% BHU PYQ -- Auto-generated & error-corrected
\documentclass[11pt,a4paper]{article}

\usepackage[a4paper,top=2.5cm,bottom=2.5cm,left=2.8cm,right=2.8cm,headheight=14pt]{geometry}
\usepackage{amsmath,amssymb}
\usepackage[T1]{fontenc}
\usepackage[utf8]{inputenc}
\usepackage{lmodern}
\usepackage{microtype}
\usepackage{fancyhdr}
\usepackage{enumitem}
\usepackage{hyperref}
\usepackage{lastpage}
\usepackage{parskip}

\hypersetup{colorlinks=true,urlcolor=black,linkcolor=black,
  pdftitle={STB-401: Sample Surveys and Design of Experiments -- BHU PYQ},pdfauthor={Banaras Hindu University}}

\pagestyle{fancy}\fancyhf{}
\renewcommand{\headrulewidth}{0.4pt}\renewcommand{\footrulewidth}{0.4pt}
\fancyhead[L]{\small   \textit{Banaras Hindu University}}
\fancyhead[R]{\small   \textit{STB-401: Sample Surveys and Design of Experi}}
\fancyfoot[C]{\small Page  \thepage\ of \pageref{LastPage}}


\newlist{parts}{enumerate}{2}
\setlist[parts,1]{label=(\alph*),leftmargin=*,itemsep=5pt,topsep=3pt}
\setlist[parts,2]{label=(\roman*),leftmargin=*,itemsep=3pt,topsep=2pt}

\newcommand{\pts}[1]{\hfill   \textit{[\#1]}}


\begin{document}


\begin{center}
  {\large\bfseries BANARAS HINDU UNIVERSITY}\\[2pt]
  {\normalsize B.Sc. (Hons.) Semester IV Examination, 2013-14}\\[6pt]
  \rule{0.8\linewidth}{0.5pt}\\[6pt]
  {\large\bfseries Paper No. STB-401: Sample Surveys and Design of Experiments}\\[4pt]
  {\normalsize Time: 3 Hours\hfill Maximum Marks: 70}
\end{center}
\rule{\linewidth}{0.4pt}
\smallskip



\noindent   \textit{Instructions: Answer   \textbf{five} questions including Question~1 which is compulsory. Symbols have their usual meanings.}
\bigskip

BS | Sem -W] 132,
Note: Answer any FIVE questions. All questions carry equal marks.
\medskip


\noindent   \textbf{Question 1.}


\begin{parts}
  \item Distinguish between sampling and complete cuumeration with examples What steps should be followed in a sample survey.
  \item A population consists of five units 10, 25, 30. 35 and 45, Give all possib'e samples of size 2 that can be drawn using SRSWR and show that the sample mean is and unbiased estimator of population mean.
\end{parts}
\medskip\rule{\linewidth}{0.2pt}

\medskip


\noindent   \textbf{Question 2.}


\begin{parts}
  \item Prove that the probability of selecting a particular unit ef the population is same for each draw in SRSWOR scheme. Also show that the sample mean square is an unbiased estimator of population mean square in SRSWOR.
  \item Describe simple random sampling of attribute. Find out the sampling variance of the unbiased estimator of population proportion in SRSWOR.
\end{parts}
\medskip\rule{\linewidth}{0.2pt}

\medskip


\noindent   \textbf{Question 3.}


\begin{parts}
  \item Define an unbiased estimator of population mean in stratified random sampling anc obtain its sampling variance under proportional allocation.
  \item Explain the purpose of stratification in sample surveys. Show that the estimator of population mean under propertional allocation provides better estimate than simple random sampling.
\end{parts}
\medskip\rule{\linewidth}{0.2pt}

\medskip


\noindent   \textbf{Question 4.}


\begin{parts}
  \item Describe systematic sampling and its utility. Compare the efficiency of systematic sampling with that of simple random sampling.
  \item What is cluster sampling?Compare the relative efficiency of cluster sampling with that of simple random sampling.
\end{parts}
\medskip\rule{\linewidth}{0.2pt}

\medskip


\noindent   \textbf{Question 5.}


\begin{parts}
  \item Derive the expression of the situation for which ratio estimator is preferred over the sample mean estimator. If coefficients of variation of study and auxiliary variables are Cy = 0.02 and C, =0.01 respectively. Suggest the appropriate estimator for estimating the population mean of study variable, where the correlation coefficient between study and auxiliary variables is p = 0.8.
  \item Describe the regression method of estimation. Find out the mean square etror of the regression estimator of population mean and compare it with that of product estimator.
\end{parts}
\medskip\rule{\linewidth}{0.2pt}

\medskip


\noindent   \textbf{Question 6.}


\begin{parts}
  \item What is analysis of variance?Describe it with a suitable example.
  \item Work out the analysis of variance for two-way classified data with one observation per cell, giving its model and complete analysis.
\end{parts}
\medskip\rule{\linewidth}{0.2pt}

\medskip


\noindent   \textbf{Question 7.}


\begin{parts}
  \item Explain the purpose of design of experiments. Describe the following fundamenta! Principles:
  \item Randomization, (i) Replication and (iii) Loca! control.
  \item Give layout and complete statistical analysis of completely randomized design. What are the advantages and disadvantages of this design?
\end{parts}
\medskip\rule{\linewidth}{0.2pt}

\medskip


\noindent   \textbf{Question 8.} Write short notes on any three of the following:


\begin{parts}
  \item Basic principles of sample survey
  \item Sampling and non-sampling errors
  \item Optimum allocation
  \item Product method of estimation
  \item Latin square design
\end{parts}

\vfill
\rule{\linewidth}{0.4pt}

\begin{center}\small
  \href{https://bhu-exam.vercel.app/}{Click here to visit our website}\\[4pt]\href{https://me-aryan.vercel.app/}{Click here to visit our contributor}\\[4pt]\href{https://quashan-me.vercel.app/}{Click here to visit our contributor}
\end{center}

\end{document}